\fchapter{Ej 7 Pág 44}
\fsection{\boldit{\textcolor{waveBlue2}{Investiga}}} 

\large{\textbf{Busca información sobre el día del sobregiro de la Tierra y responde a las siguientes cuestiones:}	 
	\begin{enumerate}[label=\alph*)]
	\item{¿En qué consiste? ¿Cuáles son sus principales causas y consecuencias?}
\begin{solution}
Día que marca la fecha en la que la humanidad a consumido todos los recursos biológicos que la Tierra puede regenerar
de manera renovable en un año. 
\vspace{1em}

Es decir, a partir de esta fecha vivimos ``a credito'' con el planeta, generando un déficit
ecológico y agotando las reservas naturales del planeta. 
\vspace{1em}

En 2025 esta fecha global cayó al 24 de julio, lo que significa que hemos usado el equivalente a 1,8 Tierras durante el año.
En caso de que toda la humanidad consumiera al ritmo actual, se necesitarían 1,8 planetas para sostenernos de forma indefinida.

\vspace{1em}
El cálculo se basa en la comparación entre la huella ecológica (la demanda humana de recursos como alimentos, madera, fibra y 
absorción de CO$_2$) y la biocapacidad de la Tierra (su capacidad para regenerar esos recursos y absorber residuos).
Principales causas:
\vspace{1em}

\phantomsection
\addcontentsline{toc}{section}{Principales causas:}
\(\boxed{\text{Principales causas:}}\)

\vspace{1em}

\begin{itemize}[label=$\bullet$]
	\item{\boldit{Sobrepoblación y patrones de consumo excesivo:} Especialmente en países ricos, donde el consumo per cápita es alto} 
	\item{\boldit{Emisiones de carbono:} La principal demanda ecológica (alrededor del 60\% de la huella global), 
		derivada de la quema de combustibles fósiles para energía, transporte e industria.}
	\item{\boldit{Pérdida de biodiversidad y degradación de suelos:} Debido a la agricultura intensiva, deforestación y urbanización, que reducen 
	la biocapacidad del planeta.}
	
\end{itemize}
 Desde 1970, esta fecha ha avanzado más de 5 meses (de finales de diciembre a julio), impulsada por
el crecimiento económico y demográfico sin avances proporcionales en eficiencia.
\end{solution}
\newpage
\vspace{1em}
\begin{solution}
\phantomsection
\addcontentsline{toc}{section}{Principales consecuencias}
\(\boxed{\text{Principales consecuencias:}}\)
\vspace{1em}

\begin{itemize}[label=$\bullet$]
	\item{\boldit{Déficit ecológico acumulado:} Equivale a 22 años de la productividad total de la Tierra; si continuamos así, crecerá 0,8 ``planetas'' por año.}
	\item{\boldit{Degradación ambiental:} Deforestación, erosión del suelo, pérdida de biodiversidad y acumulación de CO$_2$ en la atmósfera (aumento de 100 ppm desde el inicio del sobregiro global en 1971)}
	\item{\boldit{Impactos socioeconómicos:} Aumento de eventos climáticos extremos (olas de calor, inundaciones), declive en la producción de alimentos, escasez de agua y riesgos para la seguridad alimentaria.}
	
\end{itemize}

\end{solution}
\vspace{2em}

	\item{¿Cuáles son los países en los que se produce antes este día? ¿En qué países se produce más tarde? ¿Cuál crees que es el motivo? ¿Qué día se produce en España?}
\begin{solution}
Países donde se produce antes (los más adelantados, es decir, los de mayor consumo per cápita):
Estos son los que agotan su "presupuesto anual" más temprano, reflejando una huella ecológica elevada. Según los datos de 2025:
\begin{itemize}[label=$\bullet$]
	\item{Singapur: 3 de enero (necesitaría 9,4 Tierras).}
	\item{Israel: 19 de enero (5,6 Tierras).}
	\item{Baréin y Qatar: 24 de enero (alrededor de 8-9 Tierras).}
	\item{Kuwait (febrero)}
	\item{Estados Unidos (alrededor de marzo/abril, necesitando 5 Tierras).}
\end{itemize}

Países donde se produce más tarde (los más retrasados, de menor consumo), estos mantienen un equilibrio más cercano
a su biocapacidad local o tienen huellas bajas:
\begin{itemize}
	\item{Uruguay: 17 de diciembre (casi al final del año, necesitaría solo 1,1 Tierras).}
	\item{Indonesia: 17 de agosto}
	\item{Pakistán: 14 de agosto}
	\item{India: 15 de agosto(no tiene dia oficial)}
\end{itemize}

	
\end{solution}

\newpage
\begin{solution}
El motivo principal es la desigualdad en el consumo y el desarrollo económico. Los países con fechas tempranas (como Singapur o Qatar)
son economías ricas, urbanizadas y dependientes de importaciones masivas de recursos, con altos niveles de consumo energético 
(por ejemplo, aire acondicionado en climas calurosos, transporte aéreo y dietas importadas). 
\vspace{1em}

Esto genera una huella ecológica 
desproporcionada, a menudo "externalizando" el impacto a otros países mediante importaciones. En contraste, países como Uruguay 
o Indonesia tienen economías más locales, basadas en agricultura sostenible y menor industrialización, con biocapacidad alta relativa 
(bosques, tierras fértiles) y consumos per cápita bajos. Factores como el PIB per cápita, la urbanización y la dependencia de combustibles fósiles 
explican del 70 al 80\% de las variaciones.
\vspace{1em}
En España:
El Día del Sobregiro de la Tierra para España en 2025 fue el 23 de mayo (3 días más tarde que en 2024, gracias a leves mejoras 
en eficiencia energética). Si toda la humanidad viviera como en España, se necesitarían unas 2,6 Tierras para sostenernos.
\end{solution}
	\end{enumerate}}



